\section{Теоретические основы системы верификации}

В основе разрабатываемой системы верификации лежат методы формальной
спецификации и модельной проверки распределённых систем. Формальные
методы позволяют описывать систему в виде математической модели,
задающей множество допустимых состояний и переходов между ними, а
также формулировать инварианты и свойства корректности. Однако
проверка самой спецификации не гарантирует корректность конкретной
программной реализации. Для сопоставления поведения программы с её
формальной моделью используется подход, основанный на проверке трасс
исполнения по формальной спецификации.

\subsection{Формальная спецификация на языке TLA+}

TLA+ представляет собой язык формальных спецификаций, основанный на
теории множеств, логике первого порядка и темпоральной логике
действий. Он предназначен для описания поведения дискретных систем в
терминах состояний и переходов между ними \cite{lamport08}.

В TLA+ система задаётся с помощью переменных состояния. Состояние
представляет собой присваивание значений этим переменным, а шаг ---
переход от одного состояния к другому. Поведение системы задаётся как
последовательность состояний. Предикат состояния описывает свойства
отдельного состояния, а действие --- свойства перехода между двумя
соседними состояниями. Для различения значений переменных в текущем и
следующем состояниях используются обычные и штрихованные обозначения
\cite{habrias06}.

Типичная структура спецификации имеет вид

$$
Spec \triangleq Init \wedge \Box [Next]_{(v_1, v_2, \dots, v_n)}
$$

где $Init$ задаёт множество начальных состояний, $Next$ определяет
допустимые переходы системы, а $(v_1, v_2, \dots, v_n)$ --- множество
переменных состояния. Темпоральный оператор $\Box$ выражает требование
истинности формулы на всём протяжении поведения системы.

Практическая работа со спецификациями TLA+ обычно выполняется с
использованием средства проверки моделей TLC. Оно строит пространство
достижимых состояний, проверяет инварианты и позволяет обнаруживать
контрпримеры, нарушающие заданные свойства системы. В стандартном
сценарии TLA+ спецификация используется для проверки свойств самой
модели. В рамках настоящей работы спецификация рассматривается также
как эталон допустимого поведения, с которым затем сопоставляется
наблюдаемое поведение программной реализации.

\subsection{Проверка трасс исполнения по формальной модели}

Одной из важных практических задач при использовании TLA+ является
проверка соответствия программной реализации её формальной
спецификации. В идеальном случае хотелось бы доказать, что реализация
является уточнением высокоуровневой спецификации, то есть что все
возможные поведения программы допустимы спецификацией. Однако на
практике такой подход сталкивается с серьёзными проблемами
масштабируемости и применим лишь к ограниченному кругу сравнительно
небольших систем.

Полная сквозная верификация остаётся крайне трудоёмкой и технически
сложной задачей вне зависимости от используемых инструментов. Даже в
тех случаях, когда она теоретически осуществима, затраты времени и
ресурсов делают её непрактичной для большинства прикладных систем
\cite{cirstea2024validating}.

Поэтому в инженерной практике используется более прагматичный подход,
основанный на проверке трасс исполнения по формальной модели
\cite{howard2011tracechecking}. Его идея заключается не в
доказательстве того, что \emph{все} возможные поведения реализации
соответствуют спецификации, а в проверке того, что \emph{наблюдаемые}
поведения не противоречат ей.

Подход включает следующие этапы:

\begin{enumerate}
    \item программная реализация инструментируется таким образом, чтобы
    фиксировать существенные события и изменения состояния;
    \item во время выполнения системы формируется трасса исполнения;
    \item полученная трасса интерпретируется как ограничение на
    допустимое поведение формальной спецификации;
    \item средство модельной проверки определяет, существует ли
    поведение спецификации, согласованное с наблюдаемой трассой.
\end{enumerate}

Таким образом, задача сводится к проверке воспроизводимости
наблюдаемой последовательности переходов в рамках формальной модели.
В отличие от классической модельной проверки, здесь не требуется
перебирать всё пространство состояний реализации: анализируется лишь
конкретная трасса исполнения. Это существенно снижает вычислительную
сложность и делает метод применимым к распределённым системам, где
трассы могут собираться с нескольких узлов и объединяться для
последующего анализа.

Данный подход обладает рядом практических преимуществ:

\begin{itemize}
    \item относительно низкая стоимость интеграции в существующий
    процесс разработки;
    \item возможность использования стандартных средств модельной
    проверки;
    \item обнаружение некорректных переходов даже при отсутствии
    явных отказов выполнения;
    \item применимость к системам, реализованным на различных языках
    программирования.
\end{itemize}

Следует отметить, что такой метод не даёт полной формальной гарантии
корректности, поскольку анализируется лишь конечное множество
наблюдаемых трасс. Тем не менее при достаточном количестве тестовых
запусков вероятность обнаружения ошибок существенно возрастает. Как
показывает практика применения формальных методов в индустрии,
частичная формальная проверка позволяет выявлять серьёзные ошибки на
ранних этапах разработки \cite{howard2024ccf}.

\subsection{Существующие подходы}

Проверка трасс исполнения на соответствие высокоуровневым свойствам
или спецификациям имеет длительную историю в области формальных
методов. К.~Хавелунд предложил верификацию во время выполнения как
облегчённый подход, позволяющий проверять, согласуется ли трасса
исполнения системы с её высокоуровневой спецификацией
\cite{havelund2000runtime}. Как правило, такой подход предполагает
построение монитора (monitor) на основе спецификации; далее монитор получает
события трассы и проверяет их на соответствие требованиям
\cite{falcone2013runtime}.

Ховард и соавторы предложили проверять трассы исполнения
непосредственно относительно высокоуровневой спецификации системы
\cite{howard2011tracechecking}. В их работе также было показано, что
для такой проверки могут использоваться стандартные средства модельной
проверки, в частности ProB и Spin. Однако в данной работе не
рассматривались распределённые программы, для которых требуется
использование централизованных или распределённых часов для сохранения
причинно-следственных связей, а также трассы с неполной информацией.

Тасиран и соавторы одними из первых извлекали и проверяли трассы,
полученные из аппаратного симулятора, на соответствие TLA+
спецификации, тем самым продемонстрировав практическую применимость
валидации трасс \cite{tasiran2002using}. Распространение TLA+ среди
практиков распределённых систем, во многом обусловленное работой
Ньюкомба и соавторов \cite{newcombe2015aws}, а также формализация
валидации трасс как проверки уточнения, предложенная Пресслером
\cite{pressler2020conjunction}, привели к использованию данного
подхода в реальных распределённых системах.

Так, Дэвис и соавторы применили валидацию трасс к MongoDB и
обнаружили нетривиальную ошибку реализации
\cite{davis2020extreme}. При этом авторы столкнулись с трудностями
последовательного журналирования состояния реализации и согласования
различной гранулярности атомарности между спецификацией и программой.
Ниу и соавторы также использовали проверку трасс для валидации
ZooKeeper, подтверждая соответствие реализации спецификации
\cite{niu2022zookeeper}. Аналогично, Ван и соавторы выявили несколько
ошибок реализации, воспроизводя поведения TLA+ на инструментированных
реализациях \cite{wang2023model}. Их работа также демонстрирует
сложности, связанные с согласованием различной гранулярности
атомарности между спецификацией и кодом. Тем не менее все указанные
работы показывают, что сопоставление трасс реализации с
высокоуровневыми поведениями TLA+ позволяет находить нетривиальные
ошибки в реальных системах \cite{ongaro2014search}.

Для более тесного согласования высокоуровневых спецификаций с их
практическими реализациями Хэкетт и соавторы
\cite{hackett2023pgo}, а также Фу и соавторы
\cite{foo2023choreographic} предложили расширения языка PlusCal,
который используется как алгоритмическая надстройка над TLA+. Эти
расширения направлены, соответственно, на добавление в спецификации
достаточного уровня детализации для генерации кода и на генерацию
мониторов времени выполнения. Несмотря на перспективность таких
подходов, необходимость введения дополнительных прослоек между
спецификацией и реализацией может затруднять их широкое промышленное
применение. Кроме того, проекция глобальной машины состояний, которая
является распространённым стилем моделирования, на локальные машины
состояний отдельных узлов затрудняет проверку глобальных свойств с
помощью локальных мониторов.

С рассматриваемым направлением также связаны методы тестирования
реализаций по формальным спецификациям, например подходы на основе
соответствия входов и выходов \cite{tretmans1996test}. Их общая цель
также состоит в повышении уверенности в корректности реализации, однако
в таких подходах реализация обычно рассматривается как чёрный ящик,
допускающий лишь наблюдения на интерфейсах. Аналогичное замечание
относится и к методам активного обучения конечных автоматов
\cite{vaandrager2017modellearning}. В отличие от этого, при валидации
трасс предполагается наличие доступа к исходному коду реализации и
возможность его инструментирования, что даёт разработчику больше
средств для преодоления расхождений между гранулярностью атомарности
спецификации и программы.

Таким образом, существующие подходы к верификации соответствия TLA+
спецификации и исходного кода можно условно разделить на несколько
направлений: верификация во время выполнения с использованием
мониторов, проверка трасс исполнения по высокоуровневой формальной
модели, расширения языков спецификаций для генерации кода и мониторов,
а также методы тестирования и обучения моделей по наблюдаемому
поведению. Среди них наибольший интерес для настоящей работы
представляет именно проверка трасс исполнения по формальной модели,
поскольку она сочетает практическую применимость, сравнительно низкую
стоимость внедрения и возможность использовать уже существующую TLA+
спецификацию как эталон допустимого поведения реализации.

